% ---------- Carátula ----------
\begin{titlepage}
  \centering
  {\logoportada\par}
  \vspace{0.8cm}
  {\large\bfseries Universidad Abierta Interamericana\par}
  {\large Facultad de Tecnología Informática\par}
  \vspace{0.3cm}
  {\logofacultad\par}
  \vspace{1.5cm}
  {\LARGE\bfseries \nombreproyecto\par}
  \vspace{0.5cm}
  {\large \titulodoc\par}
  {\large \tipodoc\par}
  \vspace{2.5cm}
  {\large \autor\par}
  {Legajo: \legajo\par}
  \ifx\sede\empty\else{Sede: \sede\par}\fi
  \vspace{1cm}
  {Profesores: \profesores\par}
  \vfill
  {\fechadoc\par}
\end{titlepage}

% ---------- Ficha de seguimiento y revisión ----------
% Cada alumno puede definir sus filas en ficha_revision.tex (en su carpeta); si no existe, las tablas quedan vacías.
\newcommand{\filascambios}{ & & & \\ \hline}
\newcommand{\filasrevisiones}{ & & & \\ \hline}
\newcommand{\fechacreacion}{}
\IfFileExists{ficha_revision.tex}{\input{ficha_revision}}{}
\section*{Ficha de Seguimiento y Revisión}
\addcontentsline{toc}{section}{Ficha de Seguimiento y Revisión}

\medskip\noindent\textbf{Cambios del autor}\par\smallskip
\begin{tabularx}{\textwidth}{|l|X|l|X|}
\hline
\textbf{Fecha} & \textbf{Autor} & \textbf{Versión} & \textbf{Referencia de los cambios}\\ \hline
\filascambios
\end{tabularx}

\medskip\noindent\textbf{Revisiones}\par\smallskip
\begin{tabularx}{\textwidth}{|X|l|X|l|}
\hline
\textbf{Auditor} & \textbf{Versión aprobada} & \textbf{Responsabilidad/Rol} & \textbf{Fecha}\\ \hline
\filasrevisiones
\end{tabularx}

\medskip\noindent\textbf{Descripción del documento}\par\smallskip
\begin{tabularx}{\textwidth}{|l|X|}
\hline
\textbf{Ítem} & \textbf{Detalle}\\ \hline
Título del documento & \titulodoc\\ \hline
Autor & \autor\\ \hline
Fecha de creación & \fechacreacion\\ \hline
Última actualización & \fechadoc\\ \hline
\end{tabularx}

\clearpage
\tableofcontents
\clearpage

% ---------- A. Resumen ejecutivo ----------
% Máximo tres páginas (ver 00_resumen_ejecutivo/requisito.md).
\capredactadotrue
\section*{Resumen Ejecutivo}
\addcontentsline{toc}{section}{Resumen Ejecutivo}

El crecimiento de la comunidad de juegos de rol de mesa, impulsado en parte por títulos como \textit{Baldur's Gate 3}, por otros medios, como la película \textit{Dungeons \& Dragons: Honor Among Thieves}, y por producciones como \textit{Critical Role} o \textit{Dimension 20}, generó un aumento del interés de nuevos jugadores por este tipo de experiencias. Sin embargo, quienes ingresan al pasatiempo suelen enfrentarse a una curva de aprendizaje elevada, en especial al asumir el rol de \textit{Dungeon Master} (DM), por la gran cantidad de reglas, información y recursos necesarios para coordinar una campaña.

Hoy gran parte de la gestión de las partidas se realiza con herramientas dispersas, como documentos, aplicaciones de notas, manuales físicos y hojas de personaje en papel. Esto dificulta la organización y constituye una barrera de entrada para los nuevos jugadores.

Como respuesta, el proyecto propone el desarrollo de una aplicación multiplataforma orientada a la gestión integral de campañas de juegos de rol de mesa. La aplicación centraliza campañas, sesiones, personajes, inventarios, estadísticas, notas y recursos narrativos en un entorno digital accesible y organizado, y mantiene la esencia de la experiencia presencial del juego.

Entre sus funcionalidades principales se incluyen la creación y administración de campañas, la planificación de sesiones, la gestión de personajes, las notas privadas y compartidas, la personalización de recursos y el manejo de eventos en tiempo real. El proyecto busca posicionarse como una herramienta complementaria para jugadores y DMs, enfocada en la organización y la accesibilidad, y no como un reemplazo de la experiencia tradicional.

El modelo de negocio es de tipo \textit{freemium}: las funcionalidades básicas están disponibles de forma gratuita, y los ingresos provienen de una suscripción premium anual, orientada a los DMs con mayor experiencia, y, a partir de la tercera etapa del plan, de las comisiones de un \textit{marketplace} integrado. En él la comunidad ofrece campañas, mapas, recursos narrativos y otros elementos creados por los usuarios.

La inversión inicial requerida es reducida (alrededor de USD 27.610, con fondos propios) y se destina a infraestructura propia de bajo costo, registros legales y un capital de trabajo que actúa como reserva. Se trata de un modelo escalable, con potencial de crecimiento a través del contenido digital y la participación de la comunidad. Con las proyecciones a tres años se estima recuperar la inversión en alrededor de 26 meses, un plazo sujeto al nivel de adopción y de crecimiento de usuarios de la plataforma.

El proyecto es desarrollado por Tomás Bassi, estudiante avanzado de Ingeniería en Sistemas, y Manuel Regiardo, analista de sistemas y estudiante avanzado de Ingeniería en Sistemas, ambos con experiencia en desarrollo \textit{full stack}, infraestructura y práctica en juegos de rol de mesa. Ya se dispone de un prototipo funcional y de validaciones iniciales con potenciales usuarios, quienes demostraron interés en una solución enfocada específicamente en mejorar la organización y la accesibilidad de las partidas presenciales.

\clearpage
